\documentclass[conference,a4paper,10pt]{IEEEtran}

%%% ============ Packages ============
\usepackage[utf8]{inputenc}
\usepackage[T1]{fontenc}
\usepackage[spanish,es-tabla,es-nodecimaldot]{babel}
\usepackage{textcomp}
\usepackage{lmodern}
\usepackage{microtype}
\usepackage{graphicx}
\usepackage{xcolor}
\usepackage{colortbl}
\usepackage{array}
\usepackage{booktabs}
\usepackage{tabularx}
\usepackage{float}
\usepackage{tikz}
\usetikzlibrary{shapes.geometric, arrows.meta, positioning, calc, fit, backgrounds}
\usepackage[colorlinks=true,allcolors={hai-primary},urlcolor={hai-primary}]{hyperref}
\usepackage{url}
\usepackage{enumitem}
\usepackage{tcolorbox}
\tcbuselibrary{skins,breakable}
\usepackage{caption}
\usepackage{ragged2e}

%%% ============ Colors (project palette) ============
\definecolor{hai-primary}{HTML}{C2185B}
\definecolor{hai-primary-dark}{HTML}{880E4F}
\definecolor{hai-primary-soft}{HTML}{FCE4EC}
\definecolor{hai-accent}{HTML}{7B5EA7}
\definecolor{hai-accent-soft}{HTML}{F0E6F5}
\definecolor{hai-text}{HTML}{2D1B2E}
\definecolor{hai-muted}{HTML}{8B7A8E}
\definecolor{hai-warn}{HTML}{C27803}
\definecolor{hai-warn-soft}{HTML}{FEF3C7}
\definecolor{hai-rule}{HTML}{EDE0E2}
\definecolor{hai-bg}{HTML}{FAF6F7}

%%% ============ Section heading colors ============
\usepackage{titlesec}
\titleformat{\section}{\normalfont\normalsize\bfseries\scshape\color{hai-primary}\centering}{\thesection.}{0.5em}{}
\titleformat{\subsection}{\normalfont\normalsize\itshape\color{hai-text}}{\thesubsection.}{0.5em}{}
\titleformat{\subsubsection}{\normalfont\normalsize\itshape\color{hai-text}}{\thesubsubsection)}{0.5em}{}

%%% ============ Caption style ============
\captionsetup{font=small,labelfont={sc,bf,color=hai-primary},labelsep=period,justification=centering}
\captionsetup[table]{position=above,skip=4pt}
\captionsetup[figure]{position=below,skip=4pt}

%%% ============ Hyperref tweak ============
\hypersetup{
  pdftitle={HealthAI Breast BDI - Informe Tecnico},
  pdfauthor={Agustin Peralta Guarin},
  pdfsubject={Big Data, prediccion de riesgo de cancer de mama},
  pdfkeywords={Big Data, AWS, Spark, FastAPI, React, Coolify}
}

%%% ============ Custom callouts ============
\newtcolorbox{callout}[1][]{
  enhanced, breakable, sharp corners,
  colback=hai-primary-soft, colframe=hai-primary,
  boxrule=0pt, leftrule=2pt,
  left=8pt, right=8pt, top=4pt, bottom=4pt,
  fontupper=\footnotesize, #1
}
\newtcolorbox{warning}[1][]{
  enhanced, breakable, sharp corners,
  colback=hai-warn-soft, colframe=hai-warn,
  boxrule=0pt, leftrule=2pt,
  left=8pt, right=8pt, top=4pt, bottom=4pt,
  fontupper=\footnotesize, #1
}

%%% ============ 5V badge ============
\newcommand{\vtag}[1]{%
  \tikz[baseline=(b.base)]{\node[fill=hai-primary,text=white,rounded corners=4pt,inner sep=2pt,font=\scriptsize\bfseries\scshape] (b) {#1};}\hspace{4pt}%
}

%%% ============ Code style ============
\newcommand{\code}[1]{{\ttfamily\small\textcolor{hai-primary-dark}{#1}}}

%%% ============ Title ============
\title{HealthAI Breast BDI:\\Arquitectura híbrida de Big Data para la predicción de riesgo de cáncer de mama y la analítica epidemiológica}

\author{%
  \IEEEauthorblockN{Agustín Peralta Guarín \quad Natalia Moncaleano}\\
  \IEEEauthorblockA{Programa de Ingeniería de Sistemas, Facultad de Ingenierías\\
  Universidad de San Buenaventura, Bogotá, Colombia\\
  Aplicación: \url{https://bigdata.agustinynatalia.site}}%
}

\begin{document}

\maketitle

%%% ============ Abstract ============
\begin{abstract}
Este artículo presenta el diseño, implementación y despliegue de
\emph{HealthAI Breast BDI}, una aplicación web sustentada en una
arquitectura híbrida de Big Data que integra procesamiento batch sobre
Apache Spark, ingesta streaming con AWS Kinesis Firehose y un data lake
en Amazon S3, con el objetivo de entregar al usuario final una
herramienta educativa de evaluación del riesgo de cáncer de mama y un
panorama epidemiológico del país. El modelo predictivo se entrena sobre
el \emph{Breast Cancer Surveillance Consortium Risk Factor Dataset}
(1.5\,M filas agregadas, representativas de 6.7\,M mamografías). El
sistema separa el plano de datos (AWS) del plano de servicio web (VPS
administrado con Coolify), patrón que reduce el costo operativo sin
sacrificar la capacidad de procesamiento distribuido. Se discuten la
elección de la arquitectura, el pipeline de datos, la validación con
Pydantic, el modelo de clasificación servido vía FastAPI y el
cumplimiento de las cinco V del Big Data (volumen, velocidad, variedad,
veracidad y valor).
\end{abstract}

\begin{IEEEkeywords}
Big Data, cáncer de mama, modelo predictivo, Apache Spark, AWS,
streaming, Kinesis Firehose, data lake, FastAPI, React, Coolify, salud
digital.
\end{IEEEkeywords}

%%% ============================================================
\section{Introducción}
\IEEEPARstart{E}{l} cáncer de mama es el cáncer más diagnosticado en mujeres
a nivel mundial y la primera causa de muerte oncológica femenina en
Colombia~\cite{bcsc,cac}. Una porción significativa de los casos se
identifica en estadios avanzados, donde el tratamiento es más invasivo y
el pronóstico menos favorable. El tamizaje temprano y la educación en
autoexamen constituyen intervenciones costo-efectivas; sin embargo, la
información disponible para la población general se encuentra fragmentada
entre portales clínicos, datos epidemiológicos oficiales y guías médicas
dispersas.

El presente trabajo aborda esta brecha mediante el diseño de un sistema de
Big Data que integra (i) un dataset clínico masivo de mamografías
agregadas~\cite{bcsc}, (ii) reportes epidemiológicos globales y
nacionales~\cite{globocan,cac,sivigila} y (iii) un modelo de clasificación
de riesgo servido vía API. La aplicación resultante ---\emph{HealthAI
Breast BDI}--- entrega al usuario final una evaluación personalizada, un
plan de cuidado y un panorama del país en una única interfaz web
mobile-first.

El resto del documento se organiza como sigue. La Sección~\ref{sec:contexto}
describe el contexto y los objetivos. La Sección~\ref{sec:arq} presenta la
arquitectura propuesta. La Sección~\ref{sec:fuentes} detalla las fuentes de
datos. La Sección~\ref{sec:pipe} describe el pipeline. La
Sección~\ref{sec:modelo} presenta el modelo predictivo. Las
Secciones~\ref{sec:backend} y~\ref{sec:frontend} cubren el backend y el
frontend respectivamente. La Sección~\ref{sec:deploy} detalla la
infraestructura y el despliegue. La Sección~\ref{sec:5v} discute el
cumplimiento de las 5\,V. La Sección~\ref{sec:stack} resume el stack y la
Sección~\ref{sec:lim} enumera las limitaciones identificadas. La
Sección~\ref{sec:conc} concluye y delinea trabajo futuro.

%%% ============================================================
\section{Contexto y objetivos}\label{sec:contexto}

\subsection{Objetivo general}
Diseñar e implementar un sistema de Big Data que integre datos clínicos
masivos y fuentes epidemiológicas oficiales, con el fin de entregar al
usuario final una herramienta de evaluación de riesgo personalizada,
acompañada de recomendaciones preventivas y un panorama de la situación
del cáncer de mama en Colombia.

\subsection{Objetivos específicos}
\begin{itemize}[leftmargin=*,nosep]
  \item Procesar el dataset clínico BCSC ($>$6\,M mamografías) usando Apache Spark.
  \item Entrenar y servir un modelo de clasificación de riesgo a través de una API REST.
  \item Implementar un pipeline de streaming con Kinesis Firehose para captura continua.
  \item Integrar fuentes heterogéneas: GLOBOCAN, Cuenta de Alto Costo, SIVIGILA.
  \item Desplegar el sistema en producción con dominio público y HTTPS.
\end{itemize}

\subsection{Alcance}
El sistema cubre cinco módulos visibles al usuario: \emph{Mi espacio}
(dashboard personal), \emph{Mi riesgo} (evaluación tipo wizard),
\emph{Autoexamen} (guía con temporizador), \emph{Plan semanal} (checklist
de hábitos con persistencia local) y \emph{Panorama} (dashboard
epidemiológico). Se excluye la integración con sistemas de historia
clínica electrónica y la generación de diagnósticos automatizados.

%%% ============================================================
\section{Arquitectura propuesta}\label{sec:arq}
La arquitectura adopta el patrón \emph{batch + streaming sobre data
lake}, separando el plano de datos (AWS) del plano de servicio web (VPS
administrado con Coolify). Este diseño híbrido permite asignar el mejor
servicio a cada carga: AWS para el procesamiento distribuido sobre el
dataset masivo, y un VPS económico para servir la aplicación de manera
continua. La Fig.~\ref{fig:arq} ilustra el flujo de datos a través de las
cuatro etapas del sistema.

\begin{figure*}[!t]
  \centering
  \begin{tikzpicture}[
    node distance=6mm and 9mm,
    box/.style={draw=hai-rule,fill=white,rounded corners=3pt,
      minimum width=24mm,minimum height=11mm,align=center,font=\footnotesize\bfseries,
      text=hai-text},
    primary/.style={box,draw=hai-primary,fill=hai-primary-soft,text=hai-primary-dark},
    accent/.style={box,draw=hai-accent,fill=hai-accent-soft,text=hai-accent},
    stage/.style={font=\scriptsize\bfseries\scshape,text=hai-muted},
    arr/.style={-{Latex[length=2mm,width=1.4mm]},thick,hai-primary},
  ]
    %%% Row 1
    \node[box] (capt) {Formulario\\React};
    \node[primary,right=of capt] (kine) {AWS Kinesis\\Firehose};
    \node[accent,right=of kine] (spark) {EC2 Ubuntu\\+ Apache Spark};
    \node[box,right=of spark] (pkl) {modelo\_riesgo.pkl\\(S3)};

    %%% Stage labels (above)
    \node[stage,above=2mm of capt] {Captura};
    \node[stage,above=2mm of kine] {Ingesta streaming};
    \node[stage,above=2mm of spark] {Procesamiento};
    \node[stage,above=2mm of pkl] {Modelo \& servicio};

    %%% Row 2
    \node[box,below=of capt] (api) {API\\FastAPI};
    \node[primary,below=of kine] (s3) {AWS S3\\(Data Lake)};
    \node[accent,below=of spark] (pyj) {Jupyter\\\textperiodcentered\ PySpark};
    \node[box,below=of pkl] (pred) {Endpoint\\/predict};

    %%% Arrows
    \draw[arr] (capt) -- (kine);
    \draw[arr] (kine) -- (spark);
    \draw[arr] (spark) -- (pkl);
    \draw[arr] (api) -- (s3);
    \draw[arr] (s3) -- (pyj);
    \draw[arr] (pyj) -- (pred);
    \draw[arr] (capt) -- (api);
    \draw[arr] (pkl) -- (pred);
  \end{tikzpicture}
  \caption{Pipeline batch + streaming sobre data lake. Las cargas pesadas se mantienen en AWS; el servicio web se ejecuta en un VPS con Coolify.}
  \label{fig:arq}
\end{figure*}

Cada predicción que un usuario realiza desde el formulario se envía en
tiempo real al stream de Kinesis Firehose, que la deposita comprimida en
GZIP en S3 bajo el prefijo \code{raw/predictions/YYYY/MM/DD/}. Apache
Spark, ejecutándose en una instancia EC2 dedicada, lee periódicamente
esos lotes para análisis agregado y reentrenamiento del modelo.

\begin{callout}
\textbf{\itshape Justificación del diseño híbrido.} El plan inicial
contemplaba desplegar toda la aplicación sobre la EC2 que ejecuta Spark.
Tras pruebas de carga, una instancia \emph{t3.medium} no soportó
adecuadamente la combinación de Spark, Jupyter, backend y frontend
servidos en paralelo; escalar a \emph{t3.large} excedía el presupuesto
académico. La separación AWS (datos) / VPS (web) resuelve la tensión
clásica entre capacidad de procesamiento y costo de servicio continuo.
\end{callout}

%%% ============================================================
\section{Fuentes de datos}\label{sec:fuentes}
Se integraron cuatro fuentes heterogéneas. La Tabla~\ref{tab:fuentes}
resume su origen y uso dentro del sistema.

\begin{table}[!t]
\caption{Fuentes de datos integradas en el sistema.}
\label{tab:fuentes}
\centering
\renewcommand{\arraystretch}{1.15}
\arrayrulecolor{hai-primary}
\begin{tabularx}{\columnwidth}{@{}l l X@{}}
\toprule
\rowcolor{hai-primary-soft}
\textbf{\color{hai-primary-dark}Fuente} &
\textbf{\color{hai-primary-dark}Tipo} &
\textbf{\color{hai-primary-dark}Uso} \\
\midrule
BCSC~\cite{bcsc}             & Clínico EE.UU.   & Entrenamiento del modelo. $\sim$400 MB en 3 CSV; 6.7 M mamografías. \\
GLOBOCAN 2022~\cite{globocan} & Global (IARC/OMS) & Incidencia, mortalidad y prevalencia comparadas entre países. \\
Cuenta de Alto Costo~\cite{cac} & Oficial Colombia & Cáncer de mama por régimen, departamento y grupo etario. \\
SIVIGILA~\cite{sivigila}    & Vigilancia pública INS & Casos notificados con descarga abierta. \\
\bottomrule
\end{tabularx}
\arrayrulecolor{black}
\end{table}

Del BCSC se utilizan diez variables: banda etaria de 5 años, etnia,
antecedente familiar de primer grado, edad de menarca, edad del primer
parto, densidad mamaria BI-RADS, terapia hormonal de reemplazo, estado
menopáusico, IMC en bandas y biopsia mamaria previa. El dataset codifica
los valores desconocidos con el dígito \code{9}, que son tratados
explícitamente durante la fase de limpieza.

%%% ============================================================
\section{Pipeline de datos}\label{sec:pipe}
El pipeline se ejecuta en siete pasos:
\begin{enumerate}[leftmargin=*,nosep,label=\arabic*)]
  \item \emph{Ingesta del BCSC.} Descarga de los tres archivos CSV desde el portal BCSC, consolidación y limpieza inicial en \code{01\_exploracion.ipynb} con Pandas.
  \item \emph{Limpieza y transformación.} Imputación de ``desconocidos'' (código\,9), codificación ordinal de variables categóricas, verificación de la columna de pesos (\code{count}) y exportación a Parquet.
  \item \emph{Análisis exploratorio.} Notebook \code{02\_eda.ipynb}: distribución de la variable objetivo, correlaciones ponderadas, tasas por grupo etario, densidad mamaria y antecedentes familiares.
  \item \emph{Entrenamiento.} Notebook \code{03\_modelo.ipynb}: regresión logística y Random Forest con \code{class\_weight} para manejar el desbalance entre clases, evaluado con AUC-ROC, precision y recall.
  \item \emph{Procesamiento distribuido.} Notebook \code{04\_spark\_s3.ipynb}: lectura directa desde S3 con PySpark para análisis agregado sobre volúmenes que exceden la memoria RAM de una sola máquina~\cite{spark}.
  \item \emph{Servicio del modelo.} El archivo \code{.pkl} se carga al iniciar la API FastAPI~\cite{fastapi}; cada request al endpoint \code{/predict} devuelve nivel de riesgo, factores que pesan y plan de tamizaje.
  \item \emph{Captura streaming.} Cada predicción se envía a Kinesis Firehose~\cite{firehose}, que la deposita comprimida en S3 para reentrenamiento continuo.
\end{enumerate}

%%% ============================================================
\section{Modelo de predicción}\label{sec:modelo}

\subsection{Variables de entrada}
El esquema de entrada se valida con Pydantic. La Tabla~\ref{tab:vars}
enumera las variables y sus dominios.

\begin{table}[!t]
\caption{Variables de entrada del modelo (codificación BCSC).}
\label{tab:vars}
\centering
\renewcommand{\arraystretch}{1.1}
\arrayrulecolor{hai-primary}
\begin{tabularx}{\columnwidth}{@{}l l X@{}}
\toprule
\rowcolor{hai-primary-soft}
\textbf{\color{hai-primary-dark}Variable} &
\textbf{\color{hai-primary-dark}Dominio} &
\textbf{\color{hai-primary-dark}Descripción} \\
\midrule
\code{age\_group\_5\_years}     & 1--13      & Banda etaria de 5 años (1\,=\,18-29, 13\,$\geq$\,85). \\
\code{race\_eth}                & 1--9       & Etnia auto-reportada. \\
\code{first\_degree\_hx}        & 0/1/9      & Antecedente familiar de primer grado. \\
\code{age\_menarche}            & 0/1/2/9    & Edad de la primera menstruación. \\
\code{age\_first\_birth}        & 0--4, 9    & Edad del primer parto. \\
\code{BIRADS\_breast\_density}  & 1--4, 9    & Densidad mamaria BI-RADS. \\
\code{current\_hrt}             & 0/1/9      & Terapia hormonal de reemplazo actual. \\
\code{menopaus}                 & 1/2/3/9    & Estado menopáusico. \\
\code{bmi\_group}               & 1--4, 9    & IMC en bandas. \\
\code{biophx}                   & 0/1/9      & Biopsia mamaria previa. \\
\bottomrule
\end{tabularx}
\arrayrulecolor{black}
\end{table}

\subsection{Salida}
El endpoint devuelve un objeto con los campos: \code{nivel\_riesgo}
(Bajo / Medio / Alto), \code{probabilidad} (float), una lista de
\code{recomendaciones} preventivas, los \code{factores\_detalle}
ponderados, un \code{plan\_tamizaje} sugerido (frecuencia y próxima
acción) y un \code{aviso} legal.

\subsection{Métricas y validación}
El modelo se evalúa con AUC-ROC, precision y recall sobre una partición
80/20 ponderada por la columna \code{count} del dataset. Dada la
naturaleza desbalanceada del problema (la clase positiva es
minoritaria), \emph{recall} es la métrica de control: se prefieren
falsos positivos antes que falsos negativos. Pydantic rechaza con HTTP
422 cualquier valor fuera del dominio antes de que llegue al modelo, lo
que previene predicciones inválidas por error de entrada.

\begin{warning}
\textbf{\itshape Aviso clínico.} La salida del modelo es una estimación
estadística poblacional. No constituye un diagnóstico y no debe
utilizarse como reemplazo de evaluación médica profesional ni de los
protocolos de tamizaje vigentes.
\end{warning}

%%% ============================================================
\section{Backend --- API REST}\label{sec:backend}
El backend está implementado en \emph{FastAPI}~\cite{fastapi} sobre
Python\,3.10+, servido con Uvicorn y empaquetado en una imagen Docker.
Expone dos endpoints principales:
\begin{itemize}[leftmargin=*,nosep]
  \item \code{GET /health} --- verificación de estado del servicio.
  \item \code{POST /predict} --- recibe los diez factores de riesgo,
  valida con Pydantic, ejecuta el modelo cargado en memoria, registra la
  predicción en Kinesis y devuelve la respuesta estructurada al cliente.
\end{itemize}
Las llamadas a Kinesis Firehose se realizan con \code{boto3} en modo
\emph{fire-and-forget}: si el envío falla, se registra el error pero la
respuesta al usuario no se ve afectada. Esta estrategia desacopla la
captura analítica de la disponibilidad del servicio sincrónico. La
latencia del endpoint se mantiene por debajo de 200\,ms en el percentil
95 bajo carga nominal.

%%% ============================================================
\section{Frontend --- aplicación web}\label{sec:frontend}
Se construyó una \emph{Single Page Application} con React\,18 y Vite,
ruteo con React Router, microinteracciones con Framer Motion y
visualizaciones con Recharts. La aplicación se sirve a través de Nginx,
que también actúa como \emph{reverse proxy} hacia el backend para las
rutas \code{/api/*}. Los módulos visibles son:
\begin{itemize}[leftmargin=*,nosep]
  \item \emph{Inicio:} landing con introducción al sistema y CTAs.
  \item \emph{Mi riesgo:} wizard de evaluación paso a paso con persistencia local.
  \item \emph{Resultado:} nivel de riesgo, factores detalle y plan personalizado.
  \item \emph{Autoexamen:} guía con temporizador y pasos visuales.
  \item \emph{Plan semanal:} checklist de hábitos con \code{localStorage}.
  \item \emph{Panorama:} dashboard epidemiológico con descarga de datos.
  \item \emph{Proyecto:} documentación técnica navegable (página actual).
  \item \emph{Manual:} videos de uso del sistema.
\end{itemize}
La aplicación es \emph{mobile-first}: la jerarquía visual se construye
para pantallas pequeñas y escala hacia tablet y escritorio. Todos los
elementos interactivos cumplen con un área táctil mínima de
$44\times 44$ puntos y el contraste sigue las pautas WCAG.

%%% ============================================================
\section{Infraestructura y despliegue}\label{sec:deploy}

\subsection{Plano de datos --- AWS}
\begin{itemize}[leftmargin=*,nosep]
  \item \textbf{Amazon~S3} --- data lake con prefijos \code{raw/}, \code{clean/} y \code{models/}.
  \item \textbf{Kinesis~Firehose} --- ingesta DirectPUT con destino S3, buffer por tiempo/tamaño y compresión GZIP~\cite{firehose}.
  \item \textbf{EC2 (t3.medium, Ubuntu 22.04)} --- nodo de procesamiento con Apache Spark~\cite{spark} y Jupyter.
  \item \textbf{IAM} --- rol con políticas mínimas para acceso a S3 y Firehose.
\end{itemize}

\subsection{Plano de servicio --- Coolify}
Coolify~\cite{coolify} es un \emph{PaaS self-hosted} que automatiza el
ciclo de despliegue desde GitHub. Ante cada push a \code{main}, Coolify
ejecuta de forma desatendida la siguiente secuencia:
\begin{enumerate}[leftmargin=*,nosep,label=\arabic*)]
  \item Clona el repositorio en el VPS.
  \item Construye las imágenes Docker del backend y frontend con \code{docker compose build}.
  \item Orquesta los contenedores en una red interna con Traefik como reverse proxy.
  \item Gestiona HTTPS automáticamente con Let's Encrypt.
  \item Inyecta variables de entorno (credenciales AWS, buckets, streams).
\end{enumerate}
El dominio \code{bigdata.agustinynatalia.site} apunta a través de
Cloudflare al VPS. El frontend expone Nginx en el puerto 80 y enruta
\code{/api/*} al backend FastAPI en \code{backend:8000} mediante la red
Docker interna.

%%% ============================================================
\section{Cumplimiento de las 5\,V del Big Data}\label{sec:5v}
\newcommand{\vcard}[2]{%
  \begin{tcolorbox}[
    enhanced, sharp corners, breakable,
    colback=white, colframe=hai-rule,
    boxrule=0.4pt, leftrule=2pt,
    coltitle=white, colbacktitle=hai-primary,
    title={\strut\scshape\bfseries\footnotesize #1},
    fonttitle=\footnotesize,
    halign title=flush left,
    left=6pt, right=6pt, top=3pt, bottom=3pt,
    boxsep=0pt,
    attach boxed title to top left={xshift=4pt,yshift=-1pt},
    boxed title style={sharp corners,colframe=hai-primary,
      left=4pt,right=4pt,top=1pt,bottom=1pt},
    after title={\hfill},
    fontupper=\footnotesize,
    top=12pt,
  ]%
  #2%
  \end{tcolorbox}\par\vspace{2pt}%
}

\vcard{Volumen}{Más de 6\,M mamografías del dataset BCSC, almacenadas en formato Parquet columnar y procesadas con PySpark distribuido. Los datos crudos suman $\sim$400\,MB en CSV y crecen continuamente con la captura streaming.}

\vcard{Velocidad}{Pipeline de streaming con Kinesis Firehose: cada predicción del usuario llega a S3 en menos de 60\,s. La API responde con latencia $<$200\,ms en el percentil 95.}

\vcard{Variedad}{Integración de fuentes heterogéneas: dataset clínico agregado (BCSC), reportes epidemiológicos globales (GLOBOCAN), datos oficiales colombianos (Cuenta de Alto Costo) y vigilancia pública (SIVIGILA).}

\vcard{Veracidad}{Validación con Pydantic, citación explícita de fuentes oficiales, descarga abierta de los datos para auditoría por el usuario y registro inmutable de cada predicción en el data lake.}

\vcard{Valor}{Aplicación práctica con impacto educativo: evaluación personalizada, plan de cuidado, guía de autoexamen y panorama del país integrados en una sola herramienta gratuita y accesible.}

%%% ============================================================
\section{Stack tecnológico}\label{sec:stack}
La Tabla~\ref{tab:stack} sintetiza las tecnologías empleadas, organizadas por capa funcional del sistema.

\begin{table}[H]
\caption{Stack tecnológico por capa.}
\label{tab:stack}
\centering
\renewcommand{\arraystretch}{1.2}
\arrayrulecolor{hai-primary}
\begin{tabularx}{\columnwidth}{@{}>{\bfseries}l X@{}}
\toprule
\rowcolor{hai-primary-soft}
{\color{hai-primary-dark}Capa} &
\textbf{\color{hai-primary-dark}Tecnologías} \\
\midrule
Big Data \& ML     & Apache Spark (PySpark), Jupyter, Pandas, NumPy, Scikit-learn, Joblib, Parquet. \\
Cloud (AWS)        & S3, Kinesis Firehose, EC2 (t3.medium Ubuntu 22.04), IAM, AWS CLI. \\
Backend            & FastAPI, Pydantic, Uvicorn, Boto3, Docker. \\
Frontend           & React 18, Vite, React Router, Recharts, Framer Motion, Axios, Nginx. \\
Despliegue         & Coolify, Docker Compose, Traefik, Let's Encrypt, Cloudflare. \\
Control \& calidad & Git, GitHub, dotenv, ESLint, validaciones Pydantic. \\
\bottomrule
\end{tabularx}
\arrayrulecolor{black}
\end{table}

%%% ============================================================
\section{Limitaciones}\label{sec:lim}
\begin{itemize}[leftmargin=*,nosep]
  \item \emph{Sesgo geográfico del dataset.} El BCSC se construyó en EE.UU., con una composición étnica y un patrón de tamizaje distintos a los colombianos. Es necesario calibrar el modelo con datos locales antes de aplicarlo a población nacional.
  \item \emph{Modelo agregado.} El dataset utiliza filas agregadas ponderadas por \code{count}; entrenar sobre microdatos individuales requeriría acceder a una cohorte clínica.
  \item \emph{Sin factores genéticos.} Variantes BRCA1/BRCA2 no están en el dataset; el riesgo entregado es probabilístico poblacional, no familiar.
  \item \emph{Reentrenamiento manual.} Hoy el reentrenamiento se ejecuta de forma manual sobre el data lake en Jupyter.
  \item \emph{Monitoreo limitado.} Falta instrumentación de drift, distribución de inputs y feedback del usuario en producción.
\end{itemize}

%%% ============================================================
\section{Conclusiones y trabajo futuro}\label{sec:conc}
HealthAI Breast BDI demuestra que una arquitectura híbrida de Big Data
sobre la nube puede transformar un dataset clínico masivo en una
herramienta accesible al público, sosteniendo simultáneamente las cinco
V del Big Data y un caso de uso con impacto social medible. El diseño
AWS\,+\,VPS resuelve la tensión clásica entre capacidad de procesamiento
y costo de servicio continuo, y es replicable en proyectos académicos o
iniciativas con presupuesto acotado.

El sistema cumple los objetivos planteados: integra fuentes
heterogéneas, procesa $>$6\,M registros con Spark, sirve un modelo en
producción con latencia controlada, captura cada predicción para
reentrenamiento futuro y se despliega de forma automatizada en cada
push al repositorio.

El siguiente paso natural es la calibración del modelo con datos
locales y la automatización del ciclo de reentrenamiento, completando
el bucle batch\,+\,streaming que el data lake ya soporta.

%%% ============================================================
\section*{Agradecimientos}
El autor agradece a la Facultad de Ingenierías de la Universidad de San
Buenaventura, sede Bogotá, por el acompañamiento académico durante el
desarrollo del proyecto en el marco de la asignatura de Big Data 2026.

%%% ============================================================
\begin{thebibliography}{10}
\bibitem{bcsc}
Breast Cancer Surveillance Consortium, ``Risk Factor Dataset,'' 2017.
[En línea]. Disponible: \url{https://www.bcsc-research.org/data/rfdataset}

\bibitem{globocan}
International Agency for Research on Cancer, ``GLOBOCAN 2022: Global
Cancer Observatory,'' Lyon, Francia, 2022. [En línea]. Disponible:
\url{https://gco.iarc.fr}

\bibitem{cac}
Cuenta de Alto Costo, ``Situación del cáncer de mama en Colombia,''
Bogotá, Colombia, 2023. [En línea]. Disponible:
\url{https://cuentadealtocosto.org}

\bibitem{sivigila}
Instituto Nacional de Salud de Colombia, ``Sistema de Vigilancia en
Salud Pública (SIVIGILA),'' Portal de Datos Abiertos. [En línea].
Disponible: \url{https://www.datos.gov.co}

\bibitem{fastapi}
S.~Ramírez, ``FastAPI: framework de Python para construir APIs,''
Documentación oficial, 2018--2026. [En línea]. Disponible:
\url{https://fastapi.tiangolo.com}

\bibitem{spark}
The Apache Software Foundation, ``Apache Spark --- PySpark
documentation,'' v3.x, 2024. [En línea]. Disponible:
\url{https://spark.apache.org/docs/latest/api/python/}

\bibitem{firehose}
Amazon Web Services, ``Amazon Kinesis Data Firehose Developer Guide,''
2024. [En línea]. Disponible:
\url{https://docs.aws.amazon.com/firehose/}

\bibitem{coolify}
A.~Király, ``Coolify: self-hosted PaaS,'' Documentación oficial, 2024.
[En línea]. Disponible: \url{https://coolify.io}

\bibitem{sklearn}
F.~Pedregosa \emph{et~al.}, ``Scikit-learn: machine learning in
Python,'' \emph{Journal of Machine Learning Research}, vol.~12,
pp.~2825--2830, 2011.

\bibitem{birads}
American College of Radiology, \emph{Breast Imaging Reporting and Data
System (BI-RADS) Atlas}, 5th~ed., Reston, VA, USA, 2013.
\end{thebibliography}

\end{document}
